\documentclass[11pt]{article}
\usepackage[margin=1in]{geometry}
\usepackage[T1]{fontenc}
\usepackage{lmodern,microtype,amsmath,amssymb,amsthm,booktabs}
\usepackage[hidelinks]{hyperref}
\hypersetup{pdftitle={Report 139 Polynomial order for permutations containing exactly one 1432},pdfauthor={},pdfsubject={A224182, minima skeletons, a marked split, and inverse thresholds}}
\pdfinfoomitdate=1\pdftrailerid{}\pdfsuppressptexinfo=15
\setlength{\emergencystretch}{2em}
\allowdisplaybreaks[2]
\newtheorem{theorem}{Theorem}[section]
\newtheorem{lemma}[theorem]{Lemma}
\newtheorem{proposition}[theorem]{Proposition}
\newtheorem{corollary}[theorem]{Corollary}
\theoremstyle{remark}\newtheorem{remark}[theorem]{Remark}
\newcommand{\Av}{\operatorname{Av}}
\newcommand{\std}{\operatorname{std}}
\newcommand{\lm}{\operatorname{lmin}}
\newcommand{\Cat}{\operatorname{Cat}}
% Added by the ProveIt write phase (batch 105, 5 October 2026): dated notes
% marked [write], and the array package for the reading-conventions table.
\usepackage{array}
\newcommand{\file}[1]{\nolinkurl{#1}}
\newenvironment{writenote}[1][]{\par\smallskip\begingroup\small
 \noindent\textbf{[write]}\if\relax\detokenize{#1}\relax\else~\textbf{#1.}\fi\ }%
 {\par\endgroup\smallskip}
\title{Polynomial order for permutations containing exactly one 1432}
\author{Report 139}
\date{October 2, 2026}
\begin{document}
\maketitle
\begin{abstract}
Let $b_n$ count permutations with exactly one classical occurrence of $1432$.
Reversal identifies this sequence with OEIS A224182, which is defined using
$2341$. We prove $b_n=\Theta(9^n/n^3)$ by two explicit injections and a
minima-preserving Ferrers-board correspondence. The upper injection has a
computer-free eight-case proof and an exact marked inverse. The lower
injection inflates a left-to-right minimum; an entropy estimate shows that
such minima have asymptotic mean $n/3$ in the avoidance class. If $A_n$
is the number of $1234$-avoiders, the bounds give normalized lower and upper
constants $1/2187$ and $81$. They imply an inverse-threshold bracket of
unrounded asymptotic width $11/2$. An additional, explicitly computer-assisted
result characterizes the eligible marked avoiders by adjacency and four
singleton rectangles. No leading equivalent, amplitude limit, or proposed
ratio $3/80$ is proved.
\end{abstract}

\begin{writenote}[Status and trust boundary, added 5 October 2026, batch~105]
This is bundle Report~139 of an external AI-assisted research session,
filed in ProveIt as the single-source research report
\file{a224182-unique-1432-order} (provenance, the sources as the write read
them, relation to the repository, reading conventions and the collected
non-claims in Section~\ref{u1432:sec:provenance}). It is unrefereed and not
formalized; its author line reads ``Report 139'' and its PDF author field is
empty. \emph{What rests on what.} Theorem~\ref{u1432:thm:main} (the order
$b_n=\Theta(9^n/n^3)$ with the constants $1/2187$ and $81$) has a
computer-free proof; it cites Regev's strip asymptotic~\cite{Regev} and
Krattenthaler's growth-diagram theorems~\cite{Kratt}, which are not
re-proved. Theorem~\ref{u1432:thm:image} (the exact image of the split) is
computer-assisted: a proved reduction to at most eight points, then an
exhaustive finite certificate of $2{,}074$ designated cores, replayed by two
independent programs (Section~\ref{u1432:sec:image}). No asymptotic statement
rests on a computation. The write adds three remarks with proofs:
Remark~\ref{u1432:rem:lambert} (the inverse as an instance of the
transseries volume's Lambert core), Remark~\ref{u1432:rem:nonalg} (the
generating function of A224182 is not algebraic) and
Remark~\ref{u1432:rem:mrr} (the one-step bound that the source attributes
to Mansour, Rastegar and Roitershtein is not justified by their printed
argument, which gives a factor $n^2$ instead of $n$, one power of $n$ short
of the upper order), and Section~\ref{u1432:sec:further}, the open questions
with a recomputation of $b_n/(nA_n)$ for $n\le18$.
\end{writenote}
{\small\setlength{\parskip}{0pt}\tableofcontents}

\section{The result and its scope}
An occurrence of a classical pattern is a subsequence, with no adjacency
requirement, whose relative value order is the pattern. Put
\[
 b_n=\#\{\pi\in S_n:N_{1432}(\pi)=1\},\qquad
 A_n=|\Av_n(1432)|,
 \qquad M_n=\sum_{\pi\in\Av_n(1432)}\lm(\pi),
\]
where $\lm(\pi)$ is the number of left-to-right minima. Set $A_0=1$,
$M_0=0$, and $b_n=0$ for $n<4$. We prove below that
$A_n=|\Av_n(1234)|$, including equality of the full minima-skeleton
refinements. These are equal cardinalities, not equal avoidance sets.
The classical three-row strip asymptotic of Regev~\cite{Regev} is
\begin{equation}\label{u1432:eq:regev}
 A_n\sim C_0\frac{9^n}{n^4},\qquad C_0=\frac{81\sqrt3}{16\pi}.
\end{equation}

\begin{theorem}\label{u1432:thm:main}
For every $n\ge4$,
\begin{equation}\label{u1432:eq:finite}
 M_{n-3}\le b_n\le (n+2)A_{n+2}.
\end{equation}
Moreover, $M_m\sim(m/3)A_m$. Consequently,
\begin{gather}\label{u1432:eq:theta}
 b_n=\Theta(9^n n^{-3}),\qquad
 \log b_n=n\log9-3\log n+O(1),\\
 \frac1{2187}\le\liminf_{n\to\infty}\frac{b_n}{nA_n}
 \le\limsup_{n\to\infty}\frac{b_n}{nA_n}\le81.\label{u1432:eq:ratio}
\end{gather}
All assertions of this theorem have a computer-free proof.
\end{theorem}

Reversal sends $2341$ to $1432$, preserving the number of occurrences;
thus $b_n$ is A224182~\cite{OEIS}. Nakamura~\cite[Section 3.2,
Theorem 4]{Nakamura} gave a catalytic functional equation for $2341$ and
exact initial coefficients. Conway and Guttmann~\cite[Section 5.3]{CG}
conjectured the stronger equivalent $b_n\sim C_1 9^n/n^3$, with
$C_1/C_0\approx0.0375$ and the candidate $C_1/C_0=3/80$.
Theorem~\ref{u1432:thm:main} proves the polynomial order in that prediction,
but not convergence of $b_n/(nA_n)$. The proposed $3/80$ is a ratio of
amplitudes; it is not the absolute prefactor $C_1$.

The shape-Wilf ingredient is due to Backelin, West and Xin~\cite{BWX};
the directly inspected growth-diagram proof used here is
Krattenthaler~\cite[Theorems 1--3]{Kratt}. The split is adapted from
Burstein's construction for a single global $321$~\cite{Burstein}; a
supported decreasing triple requires the separate proof in
Section~\ref{u1432:sec:split}. Known fixed-occurrence exponential-growth
results are discussed in Section~\ref{u1432:sec:prior}. No worldwide priority
claim is made.

\subsection{Provenance, sources and reading conventions}\label{u1432:sec:provenance}
\begin{writenote}[Added 5 October 2026, batch~105]
\emph{Provenance.} This report is Report~139 of the session bundle of
Reports 1--243 that ProveIt's \file{docs/incoming} intake processed as
batch~105: the archive
\file{A224182_Polynomial_Growth_Order_and_Inverse_Bounds_Source.zip}
(1{,}316{,}107 bytes, wrapper directory \file{Report139/}, 10 files),
arrival commit \file{60f54ea06}, placed in \file{e85586b7c}. The text is the
delivered \file{Report139.tex} (707 lines), shipped as \file{article.tex}.
The delivered 10-page PDF and the checksum manifest \file{manifest.json} are
not shipped (the manifest was verified at placement); both are retrievable
from the arrival commit. The package names no ProveIt commit, path or
report, so the report has no pin; it carries no ``prepared for private
review'' line and no personal data. It is a single-source report, so no merge
choice was made. The write prefixed the 34 delivered labels with
\file{u1432:} and updated every reference; it added this subsection, the
notes marked [write], Remarks~\ref{u1432:rem:lambert},
\ref{u1432:rem:nonalg} and~\ref{u1432:rem:mrr},
Section~\ref{u1432:sec:further}, a table of contents and two bibliography
entries, and typeset three binomial coefficients with \verb|\binom| instead
of \verb|\choose| (identical output; this removes an \textsf{amsmath}
warning). No statement, proof or number of the source was changed: the
added remarks are the last statements of their sections and the added
displays are unnumbered, so every theorem and equation keeps its delivered
number.

\emph{The sources, as the write read them} (5~October 2026).
\begin{itemize}
\item OEIS A224182~\cite{OEIS} (revision~\#11 of September~3, 2026): ``Number
 of permutations of length $n$ containing exactly 1 occurrence of $2341$'',
 offset~1, 18 terms, through $b_{18}=2{,}055{,}119{,}789{,}289$; Nakamura
 is its author. Reversal turns $2341$ into $1432$.
\item Nakamura~\cite{Nakamura}: Section~3.2 extends the $231$ method to
 $2341$, and Theorem~4 there is the functional equation (FE2341); located.
\item Conway and Guttmann~\cite{CG}: Section~5.3 analyses the 18 OEIS terms
 of A224182 together with 26 further approximate terms of their own, finds
 the ratio $\psi^{III}_1(n)/(n\,\psi^{III}_0(n))$ (here $b_n/(nA_n)$)
 ``going to an apparent limit of $0.0375\pm0.0005$'', and writes that it is
 tempting to guess that this is $3/80$; they conclude
 $b_n\sim C_19^n/n^3$ with $C_1\approx0.0375\cdot81\sqrt3/(16\pi)$, ``and
 perhaps $C_1=3^5\sqrt3/(2^8\cdot5\pi)$ exactly'', which is $(3/80)C_0$.
 Their Section~5.5 does name A224182 for one occurrence of $1324$ (class~V),
 as Section~\ref{u1432:sec:prior} says; so does their Section~4, which
 attributes those 17 terms to Johansson and Nakamura. Both are slips of the
 paper: A224182 is the class-III sequence.
\item Mansour, Rastegar and Roitershtein~\cite{MRR}: Theorem~3.1 and its
 proof (arXiv version~1, pp.~28--29; the same text in the published author
 manuscript, PubMed Central PMC7144682). The display quoted in
 Section~\ref{u1432:sec:prior} is printed there; Remark~\ref{u1432:rem:mrr}
 explains why its printed justification gives a factor $n^2$ instead of $n$.
\item B\'ona and Burstein~\cite{BB}: Lemma~6.3 (p.~8; its rigorous proof is
 credited there to Alin Bostan), used in Remark~\ref{u1432:rem:nonalg}. Their
 non-algebraicity theorem (Theorem~6.2) concerns one copy of $k\cdots21$, $k$
 even; their generalized injection (Theorem~4.1) and non-rationality theorem
 (Theorem~5.5) concern $321\ominus\rho$, which for length four is only
 $4321$. Neither covers $1432$ or any pattern in its symmetry class
 $\{1432,2341,3214,4123\}$.
\item Krattenthaler~\cite{Kratt}: Theorems~1--3 are in Section~2 of the arXiv
 version, as Proposition~\ref{u1432:prop:refinement} uses them (the
 growth-diagram bijection, Greene-type chain readings, conjugation swapping
 the NE and SE chain statistics). The write located them and did not
 re-derive them. Burstein~\cite{Burstein} (three pages) is the short proof
 of Noonan's count of permutations with one $321$; located.
\item Not read by the write: Regev~\cite{Regev} and the journal text of
 Backelin, West and Xin~\cite{BWX} (which the source also could not
 retrieve). For Regev's constant, exact values of $A_n$ from the hook-length
 formula give $n^4A_n/(C_09^n)=0.8975$, $0.9469$, $0.9730$, $0.9864$ at
 $n=50,100,200,400$, consistent with \eqref{u1432:eq:regev} (an observation,
 not a check of the cited theorem).
\end{itemize}

\emph{What was checked at intake.} The batch-105 dossier read the manuscript
in full and found no error. It checked by hand the eight cases of
Lemma~\ref{u1432:lem:splitavoid}, the inflation interval argument and its
$\tau_k$ version, the entropy maxima $f(1/3)=\log9$ and $2\log k$ at $u=1/k$,
the constants $1/2187$, $81$ and $11/2$, the monotonicity
\eqref{u1432:eq:monotonicity} and the restriction reduction in
Theorem~\ref{u1432:thm:image}. Its own programs checked the split
exhaustively for $n\le10$ (two sizes beyond the companion), the bounds
\eqref{u1432:eq:finite} for $n\le9$, the brute-force counts $b_n$ for
$n\le10$ against the OEIS, the collision and sole-support examples and the
$106$ versus $107$ board counts. The delivered suites passed on copies
(Section~\ref{u1432:sec:prior}, note at the end).

\emph{Relation to the repository.} Before batch~105 no file of the
repository mentioned A224182, A005802 or the papers of Conway--Guttmann,
Regev, B\'ona--Burstein or Backelin--West--Xin, and the source makes no claim
about the repository, so nothing needed correcting. The report meets the
following texts.
\begin{itemize}\raggedright
\item \file{a217057-unique-pattern-occurrences} (same collection, same
 batch): exactly one $1234$ (A217057), one $1243$ (A224179) and one $12345$
 (A224248), from bundle Reports 134--138 and 140. \emph{See also}: its
 normalizer $A_n=|\Av_n(1234)|$ and Regev's constant $C_0$ are the ones here,
 and Report~136 there derives non-algebraicity of the A224179 generating
 function from the same Lemma~6.3 of~\cite{BB} that
 Remark~\ref{u1432:rem:nonalg} uses. Otherwise the methods share nothing:
 the reports there work with a tableau spine, three-row strips and unitary
 integrals and prove amplitudes; this one uses minima skeletons, Ferrers
 boards and a Burstein-type split and proves an order. Neither source cites
 the other.
\item \file{a273821-first-pattern-failure} (collection
 \file{enumerative-combinatorics}), Part~II, Section ``A worked example:
 first failure of $1432$'' (\file{fpf:lp:sec:example}): the patterns
 $\tau_r=1\oplus\delta_r$ of Corollary~\ref{u1432:cor:general}, studied there
 through a different statistic (the first failure in $123$-avoiders). No
 shared theorem.
\item \file{a047874-long-increasing-subsequences}: $A_n=|\Av_n(1234)|$ is the
 number of permutations of $[n]$ whose longest increasing subsequence has
 length at most three, a partial row sum of A047874. No shared theorem.
\item The transseries volume
 \path{Analysis/Transseries/docs/series-and-transseries/Transseries_And_Inversion/}
 (\file{transseries_and_inversion.tex}): Corollary~\ref{u1432:cor:inverse}
 is an instance of its \file{p0:thm:lambert-core} and fits the setting of
 its \file{p0:thm:staircase}\,(1) (Remark~\ref{u1432:rem:lambert}).
\end{itemize}
No Lean or Rocq file of the repository treats pattern occurrences; no
statement of this report is formalized, and its place in the collection
confers no formal status.
\end{writenote}

\begin{writenote}[What is not claimed, collected from the source]
No leading equivalent, no convergence of $b_n/(nA_n)$, no proof of the
candidate $3/80$ (which is $C_1/C_0$, not $C_1$), no D-finiteness; the two
$o(1)$ errors in the bounds are not asserted equal or effective. No matching
upper bound for $1\oplus\delta_k$, $k\ge4$. No novelty is claimed for the
upper order, and no worldwide priority claim is made; the literature and
repository check was ``scoped, not an exhaustive worldwide literature
review''. Theorem~\ref{u1432:thm:image} is computer-assisted; the main
theorem is not. The inverse bracket's width $11/2$ concerns unrounded
endpoints, ``not an exact count of admissible integers or a certified finite
numerical threshold''. No limiting density of eligible centres. The
Backelin--West--Xin full text was not retrieved; the cited Burstein paper does
not by itself prove Lemma~\ref{u1432:lem:splitavoid}. Neither the numerical
cross-checks nor the archive seal are formal verification, a digital
signature or proof of an amplitude; ``no repository publication, outside
contact, or external submission'' is part of the delivery. The write adds:
its numerical observations (Section~\ref{u1432:sec:further}) are exact
ratios of OEIS terms or floating-point diagnostics, and certify nothing about
limits.
\end{writenote}

\medskip\noindent\begin{minipage}{\textwidth}
\emph{Reading conventions} (letters the manuscript reuses, with the tempting
false readings; no symbol was renamed).

\smallskip
\begingroup\small\renewcommand{\arraystretch}{1.1}
\noindent\begin{tabular}{@{}>{\raggedright\arraybackslash}p{0.75in}>{\raggedright\arraybackslash}p{3.45in}>{\raggedright\arraybackslash}p{1.95in}@{}}
\toprule
Symbol & Meaning here (where) & Not to be read as\\
\midrule
$b_n$ & permutations of $[n]$ with exactly one $1432$, OEIS A224182 (offset
 1; A224182 itself is stated for $2341$) & the $1243$ count A224179, which
 Report~136 (filed in \file{a217057}) also writes $b_n$\\
$A_n$ & $|\Av_n(1432)|$, proved equal to $|\Av_n(1234)|$ (A005802) &
 equality of avoidance \emph{sets}; Section~\ref{u1432:sec:board} proves
 equal cardinalities only\\
$N_{1432}(\pi)$, $N(x)$ & the number of occurrences in $\pi$; the threshold
 $\min\{n\ge4:b_n\ge x\}$ (Section~\ref{u1432:sec:inverse}) & each other\\
$M_n$, $\lm$ & $M_n$ the total number of left-to-right minima over
 $\Av_n(1432)$; $\lm(\pi)$ their number in $\pi$ & the marks $M(\sigma)$ of
 Report~136 in \file{a217057}\\
$C_0$, $C_1$, $C_\pm$ & Regev's constant $81\sqrt3/(16\pi)$; Conway and
 Guttmann's conjectured amplitude of $b_n$; the proved bounds
 \eqref{u1432:eq:constants} & $C_1$ as proved; the conjectured $3/80$ is
 $C_1/C_0$, not $C_1$\\
$a,b,c,d$ & values of the unique occurrence, $a<d<c<b$, at positions
 $i<j<k<l$ (Section~\ref{u1432:sec:split}) & \\
$B$, $C$, $D$, $X$, $Y$ & the special points of the split (Lemma~\ref{u1432:lem:splitavoid});
 in the marked inverse, $C=q_K$, $D=q_{K-1}$, $B=q_{K+1}$ & the constants
 $C_\pm$, $C_0$; $D_{m,r}$\\
$k$ & a core position (Sections~\ref{u1432:sec:split},
 \ref{u1432:sec:image}); the tail length of $\tau_k$
 (Corollary~\ref{u1432:cor:general}) & each other\\
$K$ & the marked position in a split image & the tail length $k$\\
$\delta$, $\delta_k$ & a tolerance in Section~\ref{u1432:sec:entropy};
 $\delta_k$ the decreasing permutation of length $k$ & \\
$D_{m,r}$, $H$, $f$ & avoiders with $r$ minima; the binary entropy; the
 entropy exponent (Section~\ref{u1432:sec:entropy}) & \\
$\lambda$, $L$, $T_C$ & $\log9$; $\log x$; the continuous inverse
 \eqref{u1432:eq:TC} & $L$ as a length\\
(R1)--(R4) & the four forced singleton regions, equations
 \eqref{u1432:eq:R1}--\eqref{u1432:eq:R4} & \\
$\psi^{III}_r(n)$ & Conway and Guttmann's name for the number of
 permutations with $r$ occurrences of a class-III pattern; $b_n=\psi^{III}_1(n)$,
 $A_n=\psi^{III}_0(n)$ & \\
\bottomrule
\end{tabular}
\endgroup
\end{minipage}

\section{Minima skeletons and Ferrers boards}\label{u1432:sec:board}
A left-to-right minimum is an entry smaller than every entry preceding it.
Its position and value together are a point of the permutation plot, with
position increasing rightward and value upward. Fix an admissible set
$S$ of all these points, the \emph{minima skeleton}. Remove their rows and
columns. A remaining position $h$ and remaining value $y$ form an allowed
cell precisely when
\begin{equation}\label{u1432:eq:allowed}
 (s,t)\in S\text{ exists with }s<h\text{ and }t<y.
\end{equation}
In the increasing orders of the remaining positions and values, this is a
northeast-closed Ferrers board: if a cell is allowed, every remaining cell
to its right and above it is allowed.

\begin{lemma}\label{u1432:lem:skeleton}
The permutations with exact minima skeleton $S$ are in bijection with
full matchings of this board, one point in each remaining row and column.
\end{lemma}
\begin{proof}
Every nonminimum has a smaller earlier minimum and so occupies an allowed
cell. Conversely, every point of a matching has an earlier smaller
skeleton point, and is therefore not a new minimum. Nor can it invalidate
a later skeleton minimum: it is above its earlier supporting minimum,
which is itself above every subsequent skeleton minimum. Restoring $S$
therefore gives exactly the specified skeleton.
\end{proof}

A chain is \emph{board-contained} if its smallest axis-parallel bounding
rectangle consists of allowed cells. If three matching points have
positions $h_1<h_2<h_3$ and values $x>y>z$, the decreasing chain is
board-contained exactly when $(h_1,z)$ is allowed. That is exactly the
condition that some minimum can support it on the left and below.
More precisely, if $\operatorname{supp}_{\pi}(T)$ counts \emph{all}
permutation points before $h_1$ and below $z$, then
\begin{equation}\label{u1432:eq:support}
 N_{1432}(\pi)=\sum_{T\text{ a contained decreasing triple}}
                   \operatorname{supp}_{\pi}(T).
\end{equation}
Indeed every $1432$ has a nonminimum decreasing tail. Conversely an
arbitrary supporting point has an earlier-or-equal minimum still below
$z$, so its tail is contained. The sum must not count only supporting
skeleton points: ordinary points can support the same tail too.

For an increasing triple of remaining points, its first point is both
leftmost and lowest. Its cell is allowed, so its entire bounding rectangle
is allowed. Such a triple has a smaller earlier minimum and extends to
$1234$. Conversely the last three entries of any $1234$ are nonminima.
We have proved
\begin{align*}
 \pi\in\Av(1432)&\iff\text{the matching has no contained decreasing triple},\\
 \pi\in\Av(1234)&\iff\text{the matching has no increasing triple}.
\end{align*}

\begin{proposition}\label{u1432:prop:refinement}
For each fixed exact minima skeleton $S$, the numbers of $1432$-avoiders
and $1234$-avoiders realizing $S$ are equal.
\end{proposition}
\begin{proof}
Rotate the compressed board through $180$ degrees to obtain a southwest
Ferrers shape, the French convention of Krattenthaler~\cite[Section 2]{Kratt}.
Rotation acts on a matching by reverse-complement, which fixes both
patterns $123$ and $321$ and preserves rectangle containment.

Here are the precise growth-diagram facts needed from that source.
Theorem 1 bijects fillings with at most one point in each row and column
with boundary sequences of partitions; occupied rows and columns
correspond to strict successive changes. Theorem 2 reads longest
increasing and decreasing chains in southwest rectangles from the first
row and column of their partition labels. In Theorem 3, conjugating every
boundary partition swaps these two maximum-chain statistics. The
\emph{decreasing} statistic explicitly requires its bounding rectangle
to lie in the shape. Increasing chains in a southwest shape satisfy that
requirement automatically, because their northeast endpoint lies in the
shape.

Conjugation preserves every strict boundary change, hence preserves
occupancy of each row and each column, not merely the total number of
points. It therefore restricts to full matchings of our fixed board.
Summing the equality of the two chain statistics over the unrestricted
maximum and restricting the forbidden maximum to at most two equates
avoidance of the two contained triple types. The empty matching is handled
separately and trivially. Rotate back and restore the unchanged skeleton.
\end{proof}

\begin{remark}\label{u1432:rem:notoccurrence}
This is an avoidance correspondence, not an equidistribution of
occurrence counts. For the northeast board of six columns with least
allowed values $(2,1,1,1,1,1)$, the numbers of full matchings with exactly
one contained $123$ and exactly one contained $321$ are respectively
$106$ and $107$. The companion checks these counts directly. Thus a
one-occurrence theorem cannot be inferred by transferring one contained
chain through shape-Wilf equivalence.
\end{remark}

\section{An entropy bound for the number of minima}\label{u1432:sec:entropy}
Let $D_{m,r}$ count $1234$-avoiders of size $m$ with $r$ left-to-right
minima. Delete all their minima. The remainder avoids $123$: the first
point of any increasing triple in the remainder has a smaller earlier
minimum, which would extend it to $1234$. Choose the $r$ positions and
$r$ values of the minima, whose order is necessarily decreasing, and
then the standardized remainder. This is an injective encoding, although
some choices do not produce an admissible permutation. Hence
\begin{equation}\label{u1432:eq:entropybound}
 D_{m,r}\le \binom{m}{r}^{\!2}\Cat(m-r)
            \le \binom{m}{r}^{\!2}4^{m-r}.
\end{equation}
Put $H(u)=-u\log u-(1-u)\log(1-u)$, with endpoint convention
$0\log0=0$. The standard binomial entropy bound yields
\[
 D_{m,r}\le \exp\{m f(r/m)\},\qquad
 f(u)=2H(u)+(1-u)\log4.
\]
On $(0,1)$,
\[
 f'(u)=2\log\frac{1-u}{u}-\log4,
 \qquad f''(u)=-\frac2u-\frac2{1-u}<0.
\]
Thus $f$ has its unique maximum at $u=1/3$, where $f(1/3)=\log9$.
For every fixed $\delta>0$ for which the complement is nonempty,
compactness supplies $\eta_\delta>0$ such that
\[
 \sum_{|r/m-1/3|\ge\delta}D_{m,r}
 \le (m+1)9^m e^{-\eta_\delta m}.
\]
After division by \eqref{u1432:eq:regev}, the corresponding probability is
$O(m^5e^{-\eta_\delta m})$. The minima fraction converges to $1/3$ in
probability, and its boundedness in $[0,1]$ implies convergence in mean.
Proposition~\ref{u1432:prop:refinement} transfers the entire minima distribution
to $1432$-avoiders. Therefore
\begin{equation}\label{u1432:eq:minimamean}
 M_m\sim\frac m3 A_m.
\end{equation}

\section{A lower injection by inflation}\label{u1432:sec:lower}
\begin{proposition}\label{u1432:prop:lower}
For $m\ge1$, there is an injection from $1432$-avoiders of size $m$ with
a distinguished left-to-right minimum to permutations of size $m+3$
with exactly one $1432$.
\end{proposition}
\begin{proof}
If the marked minimum has value $a$, replace it by the consecutive block
\[
 (a,a+3,a+2,a+1),
\]
raising every other old value greater than $a$ by three. The four-point
block is an interval in both position and value and has pattern $1432$.

An occurrence using no block point would have existed before inflation;
one using exactly one would contract to an old occurrence. If an occurrence
uses two or three block points, these occupy consecutive positions and
consecutive ranks \emph{within the selected pattern}. The only proper
nonsingleton intervals of $1432$ are its tail pairs $43$, $32$, and its
tail triple $432$. Thus the first and smallest point of this occurrence
would lie outside the block, before it and below it. No such point exists,
because the inflated entry was a left-to-right minimum. An occurrence
using all four block points is exactly the internal one.

That unique occurrence identifies the inflated interval. Contracting it
recovers both the original avoider and its marked minimum, proving
injectivity and $M_m\le b_{m+3}$.
\end{proof}

\subsection{A fixed pattern length corollary}
The lower construction and minima concentration extend to a descending
tail of any fixed length. For $k\ge2$, put
$\tau_k=1\oplus\delta_k=(1,k+1,k,\ldots,2)$, and let $A_n^{(k)}$
and $b_n^{(k)}$ count avoidance and exactly one occurrence, respectively.

\begin{corollary}\label{u1432:cor:general}
For every fixed integer $k\ge2$,
\[
 A_n^{(k)}=|\Av_n(12\cdots(k+1))|,\qquad
 \liminf_{n\to\infty}\frac{b_n^{(k)}}{nA_n^{(k)}}
 \ge k^{-2k-1}.
\]
No matching upper bound for general $k$ is asserted.
\end{corollary}
\begin{proof}
The same fixed-skeleton board argument, forbidding contained monotone
chains of length $k$, preserves the minima distribution and proves the
avoidance equality. Let $r$ be the number of minima. Deleting them leaves
an avoider of $12\cdots k$. By transposed Robinson--Schensted, each such
permutation is a pair of standard Young tableaux with at most $k-1$ rows.
A tableau is determined by the row occupied by each of its labels, since
entries within a row increase. Encoding both tableaux by these row-label
words gives at most $(k-1)^{2(n-r)}$ possibilities, even without imposing
the same-shape and standardness conditions on the two words. Consequently
\[
 D_{n,r}^{(k)}\le \binom{n}{r}^{\!2}(k-1)^{2(n-r)}.
\]
The entropy function $2H(u)+2(1-u)\log(k-1)$ is strictly concave,
uniquely maximized at $u=1/k$, with value $2\log k$.
For this fixed $k$, Regev's strip asymptotic has the form
$A_n^{(k)}\sim C_k k^{2n}n^{-(k^2-1)/2}$, where $C_k>0$.
The proof in Section~\ref{u1432:sec:entropy} therefore gives
$M_n^{(k)}\sim(n/k)A_n^{(k)}$.

Inflate a marked minimum into the interval $\tau_k$, increasing the size
by $k$. Every proper nonsingleton interval of $\tau_k$ lies wholly in
its descending tail: an initial segment containing its first and smallest
entry and its largest entry cannot be a proper value interval. Thus an
occurrence using between two and $k$ block points would require an outside
earlier, smaller supporter, which does not exist. The arguments for zero,
one, and all $k+1$ block points are exactly as in
Proposition~\ref{u1432:prop:lower}. The unique internal occurrence recovers
the marked source, so $M_{n-k}^{(k)}\le b_n^{(k)}$. Finally
\[
 \frac{M_{n-k}^{(k)}}{nA_n^{(k)}}\longrightarrow
 \frac1k(k^2)^{-k}=k^{-2k-1}.
\]
\end{proof}

\section{The marked split and its eight-case proof}\label{u1432:sec:split}
Let the unique occurrence in $\pi$ have positions $i<j<k<l$ and values
$a<d<c<b$. Names of core points will also denote their values when no
confusion can result. Write
\[
 \pi=U\,b\,V\,c\,W\,d\,Z,
\]
where the supporting point $a$ lies in $U$. For $0<\varepsilon<1$,
form and standardize
\begin{equation}\label{u1432:eq:split}
 q=\std\bigl(U\,(c-\varepsilon)\,V\,d\,c\,b\,W\,
                      (c+\varepsilon)\,Z\bigr),
\end{equation}
marking the retained central $c$. Thus the \emph{old} $d,b$ move to the
immediate positional neighbors of $c$; the \emph{new} $c-\varepsilon$,
$c+\varepsilon$ occupy the old $b,d$ slots. For exact integer arithmetic,
multiply old coordinates by three and replace the offsets by $\pm1$.

\subsection{The forced empty regions}
Uniqueness of the core forces the following four singleton conditions:
\begin{align}
 \{\pi_t:t<j,\ \pi_t<d\}&=\{a\}, &(R1)\label{u1432:eq:R1}\\
 \{\pi_t:i<t<k,\ \pi_t>c\}&=\{b\}, &(R2)\label{u1432:eq:R2}\\
 \{\pi_t:j<t<l,\ d<\pi_t<b\}&=\{c\}, &(R3)\label{u1432:eq:R3}\\
 \{\pi_t:t>k,\ a<\pi_t<c\}&=\{d\}. &(R4)\label{u1432:eq:R4}
\end{align}
An extra point in each respective region could replace $a,b,c,d$ in the
core to give a second occurrence.

\begin{lemma}\label{u1432:lem:splitavoid}
The split permutation $q$ avoids $1432$.
\end{lemma}
\begin{proof}
It is convenient to work in unstandardized plot coordinates. The special
points are
\[
 X=(j,c-\varepsilon),\quad D=(k-\varepsilon,d),\quad
 C=(k,c),\quad B=(k+\varepsilon,b),\quad Y=(l,c+\varepsilon).
\]
Every other point is \emph{ordinary} and unchanged. In their position
order the five special points have pattern $21354$. The only inverted
special pairs are $(X,D)$ and $(B,Y)$, and $C$ makes no inversion with
another special point. Consequently a decreasing triple has at most two
special points.

Suppose a $1432$ in the output has points $(u,v,w,z)$ in position order,
with values $u<z<w<v$. Its \emph{tail} is $(v,w,z)$. If the tail contains
none of $X,D,B,Y$, it is unchanged from the original. If $u$ is ordinary,
the whole occurrence already existed and differs from the core (the old
$b,d$ are absent from the unchanged set). If $u$ is special, replace it
by $a$, which is earlier and smaller than every special point. The same
tail again gives a distinct original occurrence. Both are impossible.

For the remaining cases use transposition of the plot. The pattern
$1432$ is its own inverse, and the construction commutes with transposition:
$X$ and $D$ are interchanged, as are $B$ and $Y$, while the first and last
tail roles are interchanged. Thus for a tail with one special point it
suffices to consider $D$ and $B$ in each of the three roles. All other
tail points in these six cases are ordinary.

\medskip\noindent\textbf{Case 1. Tail $(D,w,z)$.}
Here $u<z<w<d$, so $u$ is ordinary. Since $D$ is immediately before $C$
and $w,z$ are ordinary points after $D$, they are after old position $k$.
Replacing $D$ by old $c$ gives the distinct original occurrence
$(u,c,w,z)$.

\medskip\noindent\textbf{Case 2. Tail $(v,D,z)$.}
Again $u$ is ordinary. If $v>c$, replacement of $D$ by old $c$ gives
$(u,v,c,z)$. Otherwise $d<v<c$. If $v$ is after $j$, it is between $j$
and $k$, violating (R3). If it is before $j$, then $u$ is before $j$
and below $d$, so (R1) forces $u=a$. Now $z$ is after $k$ with
$a<z<d<c$, violating (R4).

\medskip\noindent\textbf{Case 3. Tail $(v,w,D)$.}
The supporter $u$ is ordinary. Replace $D$ by old $d$ at position $l$;
this gives an original occurrence $(u,v,w,d)$. It cannot be the core:
$w$ is before $k$, whereas core $c$ is at $k$.

\medskip\noindent\textbf{Case 4. Tail $(B,w,z)$.}
If $z>a$, then $(a,b,w,z)$ using old $b$ is a distinct original
occurrence. Thus $z<a$, and $u<z<a<d$ implies that $u$ is ordinary
before $k$. If $w<c$, then $(u,c,w,z)$ is an original occurrence.
Otherwise $c<w<b$; since $w$ is after $k$, (R3) forces it to be after
$l$. The point $z$ is then also after $l$, and $(u,c,d,z)$ is an
original occurrence, since $u<z<a<d<c$.

\medskip\noindent\textbf{Case 5. Tail $(v,B,z)$.}
The ordinary point $v$ is before $k$ and above $b>c$. By (R2) it must
precede $i$, and hence $j$. Replacing $B$ by old $b$ gives the distinct
original occurrence $(u,v,b,z)$.

\medskip\noindent\textbf{Case 6. Tail $(v,w,B)$.}
The ordinary point $w$ is before $k$ and above $b>c$. By (R2) it must
precede $i$, and hence $j$. Replacing $B$ by old $b$ gives the distinct
original occurrence $(u,v,w,b)$.

\medskip
It remains to consider two special tail points. They must be either
$(X,D)$ or $(B,Y)$. If they occupy the first and last tail roles, the
intermediate ordinary point violates (R3): it is between $j,k$ with
value between $d,c$ in the first case, or between $k,l$ with value
between $c,b$ in the second. Thus the special pair occupies the first
two or last two roles. Transposition exchanges these placements within
each pair, so only the first-two placement needs checking.

\medskip\noindent\textbf{Case 7. Tail $(X,D,z)$.}
The supporter $u$ is ordinary, before $j$, with $u<z<d$. By (R1),
$u=a$. But $z$ is ordinary after $k$ and $a<z<d<c$, contrary to (R4).

\medskip\noindent\textbf{Case 8. Tail $(B,Y,z)$.}
The ordinary point $z$ is after $l$ and below $c$, because $Y$ is
immediately above $c$ among old values. If $z>a$, it violates (R4).
Hence $z<a$, and $u<z<a<d$ shows that $u$ is ordinary before $k$.
Then $(u,c,d,z)$ is a distinct original $1432$.

Every case contradicts uniqueness or a forced region condition.
\end{proof}

\subsection{The exact marked inverse}
Let the marked position in the standardized output be $K$ and its value
$C=q_K$. Let $D=q_{K-1}$ and $B=q_{K+1}$. Its adjacent values $C-1$ and
$C+1$ occur at positions $J<K-1$ and $L>K+1$, respectively, and
$D<C<B$. Delete the two positional neighbors of $K$, replace the value
at $J$ by $B$ and that at $L$ by $D$, and standardize. This recovers
$\pi$ exactly. Thus the marked split is injective into the $(n+2)$-pointed
avoiders of size $n+2$, proving the upper bound in \eqref{u1432:eq:finite}.

The mark is essential. The two distinct sources
\[
 (3,5,4,6,8,1,9,7,2),\qquad(3,8,5,1,4,2,6,9,7)
\]
both split to $(3,5,4,6,10,1,7,2,8,11,9)$, with respective
\emph{one-based} marked positions $9$ and $4$. Each source has exactly
one occurrence. This explicit collision rules out forgetting the mark
in this construction.

\section{The constants and inverse scale}\label{u1432:sec:inverse}
Combining \eqref{u1432:eq:minimamean} with Regev's equivalent gives
\[
 \frac{M_{n-3}}{nA_n}\longrightarrow\frac1{3\cdot9^3}
  =\frac1{2187},\qquad
 \frac{(n+2)A_{n+2}}{nA_n}\longrightarrow9^2=81.
\]
This proves Theorem~\ref{u1432:thm:main}. Equivalently, with
\begin{equation}\label{u1432:eq:constants}
 C_-=\frac{C_0}{2187},\qquad C_+=81C_0,
\end{equation}
one has
\[
 (C_-+o(1))\frac{9^n}{n^3}\le b_n
 \le(C_++o(1))\frac{9^n}{n^3}.
\]
The two $o(1)$ errors are not asserted to be equal or effective.

Prepending or appending a new maximum leaves the number of $1432$
occurrences unchanged: a new maximum cannot participate in that pattern
at the first or last position of the full permutation. Both maps are
injective and their images are disjoint. Therefore
\begin{equation}\label{u1432:eq:monotonicity}
 b_{n+1}\ge2b_n\qquad(n\ge4).
\end{equation}
For real $x>0$, define the finite threshold
$N(x)=\min\{n\ge4:b_n\ge x\}$. All logarithms below are natural.
Put $\lambda=\log9$, $L=\log x$, and
\begin{equation}\label{u1432:eq:TC}
 T_C(x)=\frac{L+3\log L-3\log\lambda-\log C}{\lambda}.
\end{equation}

\begin{corollary}\label{u1432:cor:inverse}
As $x\to\infty$,
\begin{equation}\label{u1432:eq:invbracket}
 \left\lceil T_{C_+}(x)+o(1)\right\rceil
 \le N(x)\le
 \left\lceil T_{C_-}(x)+o(1)\right\rceil.
\end{equation}
In particular $N(x)=(\log x+3\log\log x)/\log9+O(1)$.
The unrounded endpoints differ by $11/2+o(1)$.
\end{corollary}
\begin{proof}
For fixed $C>0$, the eventually increasing continuous function
$C e^{\lambda t}/t^3$ has inverse
$T_C(x)+O(\log L/L)$. This follows by substitution in
$\lambda t-3\log t+\log C=L$; the residual is $O(\log L/L)$ and the
derivative tends to $\lambda>0$.

The two-sided order estimate first gives
$N(x)=(L+3\log L)/\lambda+O(1)$. Uniformly on any fixed enlargement of
this bounded-width index window, the relative errors in
\eqref{u1432:eq:constants} tend to zero. Bound their absolute values by one
positive number $\epsilon(x)\to0$, chosen also to cover the two
continuous inverse endpoints. Let $t_+$ and $t_-$ be the real crossings of $x$ by the upper and
lower envelopes $(C_++\epsilon(x))e^{\lambda t}/t^3$ and
$(C_--\epsilon(x))e^{\lambda t}/t^3$, respectively. At the integer
$\lceil t_+\rceil-1$, the upper envelope is strictly below $x$, so
$b_{\lceil t_+\rceil-1}<x$. At $\lceil t_-\rceil$, the lower
envelope is at least $x$, so $b_{\lceil t_-\rceil}\ge x$. These two
integers lie in the chosen window for large $x$. Monotonicity gives
$\lceil t_+\rceil\le N(x)\le\lceil t_-\rceil$. The inverse formula just established and
$\log(C_\pm\pm\epsilon(x))=\log C_\pm+o(1)$ give
\eqref{u1432:eq:invbracket}.

Finally
\[
 T_{C_-}(x)-T_{C_+}(x)=\frac{\log(C_+/C_-)}{\log9}
 =\frac{\log(81\cdot2187)}{\log9}=\frac{11}{2}.
\]
\end{proof}
The endpoint errors in \eqref{u1432:eq:invbracket} occur \emph{inside} the
ceilings and may differ. They cannot be discarded near an integer.
The width $11/2$ describes real endpoints, not an exact count of
admissible integers or a certified finite numerical threshold.

\begin{remark}[{[write]} The inverse in the terms of the transseries volume, 5~October 2026]\label{u1432:rem:lambert}
The equation $\lambda t-3\log t+\log C=L$ inverted in the proof of
Corollary~\ref{u1432:cor:inverse} is the dominant block
$\Phi_0(X)=aX+b\log X=L'$ of \file{p0:thm:lambert-core} in the transseries
volume~\cite{TSvol}, with
\[
 a=\lambda=\log9,\qquad b=-3,\qquad L'=L-\log C .
\]
Since $b<0$, part~(3) of that theorem says that for
$L'>L_c=3\bigl(1-\log(3/\lambda)\bigr)$ there are exactly two positive
solutions, and the large one, on which $\Phi_0$ increases, is
\[
 X=\frac3\lambda\Bigl(-W_{-1}\bigl(-\tfrac\lambda3e^{-L'/3}\bigr)\Bigr).
\]
With $u=\frac\lambda3e^{-L'/3}\to0^+$, the expansion
$-W_{-1}(-u)=\ell_1+\ell_2+O(\ell_2/\ell_1)$, $\ell_1=\log(1/u)=L'/3+\log(3/\lambda)$,
$\ell_2=\log\ell_1=\log L-\log3+O(1/L)$, gives
\[
 X=\frac{L'+3\log(3/\lambda)+3\log L-3\log3}{\lambda}+O\Bigl(\frac{\log L}{L}\Bigr)
  =T_C(x)+O\Bigl(\frac{\log L}{L}\Bigr),
\]
the formula used in the proof, so $T_C(x)$ is the two-term expansion of the
Lambert solution. The threshold $N(x)$ is the integer staircase $N_\ast$ of
\file{p0:def:three-inverses} with $n_1=4$: by \eqref{u1432:eq:monotonicity}
and $b_4=1$, $(b_n)_{n\ge4}$ is strictly increasing and unbounded, so it has
admissible interpolations (piecewise linear, for example), and
\file{p0:thm:staircase}\,(1) gives $N(x)=\lceil\nu(x)\rceil$ for each of them.
The proof of Corollary~\ref{u1432:cor:inverse} does not choose an
interpolation: it compares $b_n$ with two continuous envelopes at the
integers, which yields the two-ceiling bracket directly. No novelty is
claimed for the inversion.
\end{remark}

\begin{remark}[{[write]} The generating function is not algebraic, 5~October 2026]\label{u1432:rem:nonalg}
The generating function $B(z)=\sum_{n\ge1}b_nz^n$ of A224182 is not an
algebraic power series.

\emph{Proof.} Put $f(z)=z+z^2+z^3+B(z)$, so that $f_n=1$ for $n\le3$ and
$f_n=b_n$ for $n\ge4$. For $n\ge4$, \eqref{u1432:eq:finite} gives
$b_n\ge M_{n-3}\ge A_{n-3}\ge1$, because every permutation has at least one
left-to-right minimum; hence $f_n>0$ for every $n\ge1$. By
\eqref{u1432:eq:theta} there are $0<c'\le C'$ and $n_0$ with
$c'9^n/n^3\le f_n\le C'9^n/n^3$ for $n\ge n_0$. Put
$c=\min\{c',\min_{1\le n<n_0}f_nn^3/9^n\}>0$ and
$C=\max\{C',\max_{1\le n<n_0}f_nn^3/9^n\}$. Then
\[
 c\,\frac{9^n}{n^3}\le f_n\le C\,\frac{9^n}{n^3}\qquad(n\ge1).
\]
The series $f$ has nonnegative coefficients and radius of convergence $1/9$,
so it is analytic at the origin. Lemma~6.3 of B\'ona and
Burstein~\cite{BB}, with $K=9$ and the integer $m=3>1$, shows that $f$ is not
an algebraic power series. As $B=f-z-z^2-z^3$ differs from $f$ by a
polynomial, $B$ is not algebraic either. \qed

The source does not state this consequence of its theorem. Report~136 of the
same bundle (in \file{a217057-unique-pattern-occurrences}) draws it in the
same way for one $1243$; B\'ona and Burstein's own non-algebraicity theorem
concerns one copy of $k\cdots21$ with $k$ even and does not cover $1432$.
Whether $B$ is D-finite is open (Question~\ref{u1432:q:dfinite}).
\end{remark}

\section{An exact eligible-center statistic}\label{u1432:sec:image}
The marked inverse suggests a finite geometric statistic on avoiders.
The following refinement is logically separate from the computer-free
order theorem. Its sufficiency proof uses the exhaustive certificate
specified below.

Given $q\in\Av_{n+2}(1432)$ with a marked position $K$, require
\begin{equation}\label{u1432:eq:adjacency}
 1<K<n+2,\quad D=q_{K-1}<C=q_K<B=q_{K+1},\quad
 \operatorname{pos}_q(C-1)<K-1,\quad
 \operatorname{pos}_q(C+1)>K+1.
\end{equation}
The last conditions include existence of the adjacent values. Unsplit
as in Section~\ref{u1432:sec:split}. The result $p$ has distinguished tail
$b,c,d$ at positions $j<k<l$. Require exactly one $a<d$ before $j$,
and let its position be $i$. Finally require (R2)--(R4) in $p$.
Together with sole support (R1), these are four singleton rectangles.
Call such a center \emph{eligible}.

\begin{theorem}[Exact image with a finite certificate]\label{u1432:thm:image}
The split is a bijection between permutations of size $n$ with one
$1432$ and eligible marked centers of $1432$-avoiders of size $n+2$.
In particular,
\[
 b_n=\sum_{q\in\Av_{n+2}(1432)}
             \#\{K:K\text{ is eligible in }q\}.
\]
\end{theorem}
\begin{proof}
Necessity follows from the split and (R1)--(R4). Conditions
\eqref{u1432:eq:adjacency} ensure that all roles are distinct. Unsplit then
split again: deleting $D,B$ and replacing the $C-1,C+1$ entries removes
exactly the two new ranks adjacent to $C$, while deleting their neighbors
removes exactly the two new positions adjacent to $K$. After
standardization the forward construction therefore recovers the given
$q$ and its mark. In particular $p$ has the designated occurrence
$(i,j,k,l)$.

Suppose $p$ has another occurrence. Retain the union of this occurrence
and the designated core, at most eight points, and standardize.
All four singleton-region conditions persist: their mandatory core points
remain, and deletion introduces no extra points. Splitting commutes with
this restriction. To see this, give $X,B$ ancestor $b$, give $D,Y$
ancestor $d$, and give every other split point its own old ancestor.
Restriction to a set containing the core keeps precisely the output
points with ancestors in that set. The two new values remain immediately
adjacent to $c$ among old values; the two new positions remain immediately
adjacent to $k$ among old positions. Every comparison is preserved.
Thus splitting the restricted source and standardizing equals restricting
the split output and standardizing, and is still an avoider.

It remains to check the following finite, completely specified statement:
for each $4\le m\le8$, every permutation $p\in S_m$ having at least two
occurrences, and every designated $1432$ core passing (R1)--(R4), has a
split containing $1432$. The companion enumerates these objects in two
independently written implementations, one using word substitution and
one using integer plot coordinates. Both give the following counts;
all designated cores counted in the middle row have a nonavoiding split.
The sealed data includes every one of these $2{,}074$ designated cores,
its exact split, and one explicit forbidden quadruple in that split;
replay independently enumerates the complete list again.
\begin{center}
\begin{tabular}{lrrrrr}\toprule
Size $m$&4&5&6&7&8\\\midrule
Cores in multiple-occurrence sources&0&0&6&130&1938\\
Cores with an avoiding split&0&0&0&0&0\\\bottomrule
\end{tabular}
\end{center}
This finite certificate contradicts the reduced counterexample and
proves sufficiency in every size. The inverse already proved gives
uniqueness.
\end{proof}

Sole support does not suffice without the other three regions. For
$q=(1,3,6,2,4,7,5)$ with one-based center $5$, the inverse is
$p=(1,5,4,3,2)$, which has four occurrences. The exact eligible-center
statistic provides a possible target for amplitude analysis, but no
limiting density for this statistic is established here.

\section{Prior results and reproducibility}\label{u1432:sec:prior}
The exponential growth rate $9$ is inherited from avoidance and the
general fixed-occurrence theory of Mansour, Rastegar and
Roitershtein~\cite{MRR}. Their proof of Theorem 3.1 displays the stronger
one-step estimate $f_m^\xi(n)\le n\sum_{r=0}^{m-1}f_r^\xi(n-1)$,
where $f_r^\xi(n)$ counts exactly $r$ occurrences of $\xi$. For $m=1$
and $\xi=1432$, this would give $b_n\le nA_{n-1}$ and hence the same
upper order. We do not rely on that estimate here: the marked split
independently proves the upper bound actually used in
Theorem~\ref{u1432:thm:main}. No novelty is claimed for the upper order. The contribution established by the argument here
is the complete stated order theorem with a matched lower scale and
explicit constants, relative to the conjecture being addressed.
A scoped literature and repository check is not a global priority proof.

\begin{remark}[{[write]} The one-step bound of Mansour, Rastegar and Roitershtein, 5~October 2026]\label{u1432:rem:mrr}
The display quoted above is printed in the proof of Theorem~3.1
of~\cite{MRR} (arXiv version~1, p.~29; the same text in the published author
manuscript), where it is derived from one sentence: removing the leftmost
letter of the leftmost occurrence of $\xi$ in $\pi\in S_n$ and renaming the
remaining letters gives $\pi'\in S_{n-1}$ with at most $m-1$ occurrences.
That map recovers $\pi$ from $\pi'$ only together with the position \emph{and
the value} of the removed letter, so the printed argument bounds each fibre
by $n^2$, not by $n$, and proves
\[
 f_m^\xi(n)\le n^2\sum_{r=0}^{m-1}f_r^\xi(n-1).
\]
This suffices for Theorem~3.1 there, which concerns exponential rates only.
For $\xi=1432$ and $m=1$ it gives $b_n\le n^2A_{n-1}=O(9^nn^{-2})$, one
power of $n$ short of the upper order in~\eqref{u1432:eq:theta}.

The fibres do exceed $n$. For $n=8$, let $v\in\{1,2,3\}$ and
$p\in\{1,2,3\}$, and let $\pi$ be obtained from $4571263\in\Av_7(1432)$ by
raising every letter $\ge v$ by one and inserting $v$ at position $p$; the
nine permutations are
\[
 15682374,\ 25681374,\ 35681274,\ 51682374,\ 52681374,\ 53681274,\
 56182374,\ 56281374,\ 56381274 .
\]
A $1432$ is a letter followed by a decreasing triple of larger letters. In
each $\pi$ the letters other than $v$ read $5,6,8,x,y,7,4$ in this order,
where $\{x,y\}=\{1,2,3\}\setminus\{v\}$ and $x<y$. No letter other than $v$
starts a $1432$: the larger letters following $5$, $6$, $x$ and $y$ are
among $6,8,7$, among $8,7$, the letters $y,7,4$ and the letters $7,4$,
respectively, and $8$, $7$ and $4$ are followed by no larger letter; none of
these lists contains a decreasing triple. The larger letters following $v$
are, in order, some of $5,6$, then $8$, then some of $x,y$, then $7,4$; as
$x<y<4$, their only decreasing triple is $(8,7,4)$. Hence $\pi$ has exactly
one $1432$, namely $(v,8,7,4)$; its leftmost letter is $v$, and removing it
gives back $4571263$. So this
avoider of size $7$ has at least nine preimages in $S_8$. By exhaustive
enumeration the largest fibre is $2,4,6,9,12$ for $n=5,\dots,9$.

Whether the stated inequality with the factor $n$ holds is therefore not
settled by~\cite{MRR}; it is Question~\ref{u1432:q:mrr}. The source's
sentence that the estimate ``would give'' the same upper order is correct as
a conditional, and its non-claim of novelty for the upper order is kept as the
source's own. But among the sources the manuscript cites, the marked split of
Section~\ref{u1432:sec:split} is the only complete proof of
$b_n=O(9^n/n^3)$ that the write found. (The placement record of this batch,
commit \file{e85586b7c}, says that the upper order was known through this
paper; this remark qualifies that statement.)
\end{remark}

The original Backelin--West--Xin journal full text was not successfully
retrieved in this investigation. Attribution to that work is distinct
from the directly checked proof in Krattenthaler's Theorems 1--3, whose
hypotheses and rectangle convention were verified above. The cited
Burstein paper concerns a unique global decreasing triple; it does not
by itself prove Lemma~\ref{u1432:lem:splitavoid}. Conway--Guttmann Section 5.5
repeats identifier A224182 for a different class; the identification here
uses OEIS's $2341$ definition and reversal, consistent with their
Section 5.3 and Nakamura's equation.

\subsection{What the companion checks}
The offline archive contains this editable \LaTeX\ article and PDF,
two standard-library-only exact implementations, closed typed data,
a read-only integrity/replay/build tool, and source provenance. It needs
no network access or third-party Python package. Normal and optimized
isolated Python runs must agree byte for byte; assertions removable by
\texttt{-O} are not used as checks.

Through size eight the replay checks every permutation and all $5{,}102$
one-occurrence sources, the marked inverses and injectivity, and the
exact equality for each minima position-and-value skeleton. It also
checks $75{,}419$ core-preserving restriction commutations, all $110{,}459$
interior-center candidates in avoiders of sizes six through eight, all $314$
marked-minimum inflations through source size five, the explicit
unmarked collision, the $106$ versus $107$ board obstruction, and the
finite exact-image certificate. Mathematical fixture values are
recomputed rather than merely trusted after a hash comparison.

The main theorem does not depend on these finite tests: its split proof,
inflation argument, entropy estimate, and uses of cited avoidance
theorems are ordinary mathematical proofs. Theorem~\ref{u1432:thm:image}
\emph{does} use its finite exhaustive certificate, after a proved
universal obstruction reduction. This distinction is essential. Neither
the numerical cross-checks nor the archive seal constitute formal
verification, a digital signature, or proof of an asymptotic amplitude.

The companion's documented commands reject nonisolated launches before
nonbuilt-in imports, enforce a closed inventory and strict JSON schemas,
reject unsafe output targets and symlink paths, compare exact replay in
normal and optimized modes, and test deliberately corrupted copies.
PDF reproduction performs two fresh clean builds and compares their bytes
with the frozen PDF. Deterministic ZIP packing is checked again after
fresh extraction. These are reproducibility and integrity controls,
not a substitute for reviewing the mathematics.

\begin{writenote}[The companion here, and the intake reruns, 5~October 2026]
``The offline archive'' is the delivered package; in this report its four
programs are \file{code/verify.py}, \file{code/primary.py},
\file{code/coordinate.py} and \file{code/replay.py}, its sealed data
\file{checks/fixtures.json} is \file{data/checks-fixtures.json}
(864{,}678 bytes, the $2{,}074$ certificate records and the other fixture
values), and its source provenance is \file{SOURCES.md}; the PDF and the
manifest are not shipped. The verifier enforces the delivered layout (it
reads \file{manifest.json}, \file{checks/} and \file{Report139.pdf}), so it
does not run in the shipped directory; the report README gives two routes,
from the arrival commit's archive and from the shipped programs. At
placement, on copies, \texttt{check}, \texttt{replay} and \texttt{-O replay}
passed (43--49~s under load; both outputs byte-identical to the sealed
fixture) and \texttt{selftest} passed 51 tests in 238~s. At the write, on a
fresh extraction, \texttt{check} and \texttt{replay} (normal and \texttt{-O})
passed again in about 5~s each, with output byte-identical to the shipped
fixture, and running \texttt{replay.run()} from copies of the three shipped
mathematical programs reproduced the shipped fixture's content exactly.
The PDF build and packing commands, which compare TeX environments, were
not run.
\end{writenote}

\section{Further questions}\label{u1432:sec:questions}
The principal unresolved question is whether $b_n/(nA_n)$ has a limit,
and if so whether that limit is the suggested $3/80$. The exact image
formula in Theorem~\ref{u1432:thm:image} turns this into a question about the
asymptotic mean number of eligible centers in a uniform avoider. Even a
proof of convergence without evaluation would strengthen the order theorem.

The constants $1/2187$ and $81$ arise from transparent injections and are
not asserted to be sharp. Better control of the split image could narrow
them and the inverse bracket. For $1\oplus\delta_k$, a matching general
upper bound would complement Corollary~\ref{u1432:cor:general}; none follows
from the present eight-case proof. Finally, a computer-free sufficiency
proof for the four-rectangle exact image criterion would remove the finite
certificate from that additional theorem, just as the eight-case argument
already does for the main upper bound.

\subsection{Further questions and research}\label{u1432:sec:further}
\begin{writenote}[Added 5 October 2026, batch~105]
Following Vladimir Reshetnikov's standing rule of 4~October 2026, every
claim of the source that is not proved in full is stated here as an open
question, with its source, the argument offered and what is missing. The
three paragraphs above are
Questions~\ref{u1432:q:limit}, \ref{u1432:q:density} (first paragraph),
\ref{u1432:q:constants}, \ref{u1432:q:tau} (second) and \ref{u1432:q:image}
(third); Questions~\ref{u1432:q:mrr} and~\ref{u1432:q:dfinite} come from
Remarks~\ref{u1432:rem:mrr} and~\ref{u1432:rem:nonalg};
Question~\ref{u1432:q:external} concerns the external inputs. The intake
found no false claim in the source.
\end{writenote}
\begin{enumerate}
\item\label{u1432:q:limit} \textbf{Does $b_n/(nA_n)$ converge, and to
 $3/80$?} (Source, first paragraph; Conway and Guttmann~\cite[Section~5.3]{CG}.)
 Theorem~\ref{u1432:thm:main} confines the ratio to $[1/2187,81]$ in the
 limit; nothing in the source bears on convergence. Missing: any limit
 theorem; Question~\ref{u1432:q:density} is one route.

 \emph{Observation, not a claim.} From the 18 terms of A224182 and exact
 $A_n$ (hook-length formula), the ratios $r_n=b_n/(nA_n)$, to seven digits,
 are
 \begin{center}\small\setlength{\tabcolsep}{3.5pt}
 \begin{tabular}{@{}lccccccccc@{}}\toprule
 $n$ & 10 & 11 & 12 & 13 & 14 & 15 & 16 & 17 & 18\\\midrule
 $r_n$ & .0369735 & .0375056 & .0378384 & .0380455 & .0381713 & .0382435 & .0382801 & .0382927 & .0382892\\
 \bottomrule\end{tabular}
 \end{center}
 They increase through $n=17$ and decrease for the first time at $n=18$
 ($r_{18}-r_{17}=-3.5\cdot10^{-6}$); $r_n$ first exceeds $3/80=0.0375$ at
 $n=11$. Polynomials in $1/n$ of degrees $1,2,3,4$ through the last
 $2,3,4,5$ terms extrapolate to $0.0382$, $0.0361$, $0.0367$, $0.0370$; for
 degrees $2$--$4$ the extrapolated value increases as the last term used
 moves from $n=14$ to $n=18$. All of this
 is consistent with Conway and Guttmann's $0.0375\pm0.0005$, which used 26
 further approximate terms. The placement record of this batch
 (\file{e85586b7c}) and its dossier read the same data as ``above $3/80$ and
 rising''; the turn at $n=18$ corrects that reading. Eighteen terms cannot
 separate $3/80$ from nearby values, and none of these numbers is a bound.
\item\label{u1432:q:density} \textbf{A limiting density of eligible
 centres} (source, first paragraph and after Theorem~\ref{u1432:thm:image}).
 By Theorem~\ref{u1432:thm:image}, $b_n=A_{n+2}E_{n+2}$, with $E_{n+2}$ the
 mean number of eligible centres of a uniform $1432$-avoider of size $n+2$.
 Since $A_{n+2}/A_n\to81$, the ratio $b_n/(nA_n)$ converges if and only if
 $E_{n+2}/n$ does, and the value $3/80$ corresponds to $E_{n+2}\sim
 n/2160$. Missing: any law of large numbers for eligible centres.
\item\label{u1432:q:constants} \textbf{Sharper constants and a narrower
 inverse bracket} (source, second paragraph). The constants $1/2187$ and $81$
 come from the two injections and are not claimed sharp; the unrounded width
 $11/2$ of Corollary~\ref{u1432:cor:inverse} is $\log(81\cdot2187)/\log9$.
 If the inequality of Question~\ref{u1432:q:mrr} holds, the upper constant
 improves from $81$ to $1/9$ ($C_+=C_0/9$) and the width to
 $\log(2187/9)/\log9=5/2$.
 Missing: control of the split image (the eligible-centre statistic) or a
 better lower injection.
\item\label{u1432:q:mrr} \textbf{The one-step inequality with the factor
 $n$} (from Remark~\ref{u1432:rem:mrr}). Is $f_m^\xi(n)\le
 n\sum_{r<m}f_r^\xi(n-1)$ true, in general or for $\xi=1432$, $m=1$, that
 is $b_n\le nA_{n-1}$? The printed argument of~\cite{MRR} gives $n^2$, and
 its fibres exceed $n$ from $n=8$. \emph{Observation:} $b_n/(nA_{n-1})$
 increases from $0.0417$ at $n=4$ to $0.2787$ at $n=18$; brute force over
 every $m$, $n\le7$, for the patterns $123$, $132$, $1234$, $1243$, $1324$,
 $1432$ and $2413$ finds no violation (largest ratio $0.46$). If Conway and
 Guttmann's estimate is right, $b_n/(nA_{n-1})\to9\cdot3/80=0.3375$, inside the
 bound. Missing: a proof; counting the fibres of the deletion map cannot
 give it.
\item\label{u1432:q:tau} \textbf{A matching upper bound for
 $\tau_k=1\oplus\delta_k$, $k\ge4$} (source, second paragraph;
 Corollary~\ref{u1432:cor:general}). The lower bound
 $\liminf b^{(k)}_n/(nA^{(k)}_n)\ge k^{-2k-1}$ is proved; no upper bound of
 order $nA^{(k)}_n$ is. The eight-case proof of
 Lemma~\ref{u1432:lem:splitavoid} is specific to $k=3$, and the printed
 argument of~\cite{MRR} gives only $b^{(k)}_n\le n^2A^{(k)}_{n-1}$, a factor
 $n$ too large. Missing: a split for a longer decreasing tail, with its case
 analysis.
\item\label{u1432:q:image} \textbf{A computer-free proof of sufficiency in
 Theorem~\ref{u1432:thm:image}} (source, third paragraph). The reduction to
 at most eight points is proved; the remaining statement is the finite
 certificate of $0+0+6+130+1938=2{,}074$ designated cores in sizes $4$--$8$.
 Missing: a case analysis, like Section~\ref{u1432:sec:split}'s, showing
 that a second occurrence always survives the split.
\item\label{u1432:q:dfinite} \textbf{Is the generating function of A224182
 D-finite?} (The delivered README: the package ``does not prove \dots\
 D-finiteness''; Remark~\ref{u1432:rem:nonalg}.) It is not algebraic. The integer
 exponent $-3$ does not obstruct D-finiteness, and Nakamura's functional
 equation (FE2341)~\cite{Nakamura} uses $n(n+1)$ catalytic variables at
 size $n$.
 Missing: an annihilating operator, or an argument excluding one.
\item\label{u1432:q:external} \textbf{External inputs and priority.} The
 proofs use Regev's strip asymptotic~\cite{Regev} and Krattenthaler's
 Theorems~1--3~\cite{Kratt}, which the write located but did not re-derive,
 and attribute the shape-Wilf ingredient to Backelin, West and
 Xin~\cite{BWX}, whose journal text neither the source nor the write read.
 The source's literature check was scoped; priority for the matched order and
 the constants is not established beyond it.
\end{enumerate}

\begin{thebibliography}{99}
\small\raggedright
\setlength{\itemsep}{3pt}
\bibitem{OEIS} OEIS Foundation, \emph{A224182: Number of permutations of
length $n$ with exactly one occurrence of the pattern $2341$}.
\url{https://oeis.org/A224182}. Definition and initial terms checked
October 2, 2026.
\bibitem{Nakamura} B. Nakamura, \emph{Approaches for enumerating
permutations with a prescribed number of occurrences of patterns},
arXiv:1301.5080 (2013), Section 3.2, Theorem 4.
\url{https://arxiv.org/abs/1301.5080}.
\bibitem{CG} A. R. Conway and A. J. Guttmann, \emph{Counting occurrences
of patterns in permutations}, Electronic Journal of Combinatorics
\textbf{32}(1) (2025), P1.3, Section 5.3.
\url{https://www.combinatorics.org/ojs/index.php/eljc/article/view/v32i1p3}.
\bibitem{Burstein} A. Burstein, \emph{A short proof for the number of
permutations containing pattern $321$ exactly once}, Electronic Journal of
Combinatorics \textbf{18}(2) (2011), P21.
\url{https://www.kurims.kyoto-u.ac.jp/EMIS/journals/EJC/Volume_18/PDF/v18i2p21.pdf}.
\bibitem{BWX} J. Backelin, J. West and G. Xin, \emph{Wilf-equivalence
for singleton classes}, Advances in Applied Mathematics \textbf{38}
(2007), 133--148, Proposition 2.2.
\url{https://doi.org/10.1016/j.aam.2004.11.006}.
\bibitem{Kratt} C. Krattenthaler, \emph{Growth diagrams, and increasing
and decreasing chains in fillings of Ferrers shapes}, Advances in Applied
Mathematics \textbf{37} (2006), 404--431, Theorems 1--3.
\url{https://arxiv.org/abs/math/0510676}.
\bibitem{Regev} A. Regev, \emph{Asymptotic values for degrees associated
with strips of Young diagrams}, Advances in Mathematics \textbf{41}
(1981), 115--136.
\url{https://doi.org/10.1016/0001-8708(81)90012-8}.
\bibitem{MRR} T. Mansour, R. Rastegar and A. Roitershtein,
\emph{Finite automata, probabilistic method, and occurrence enumeration
of a pattern in words and permutations}, SIAM Journal on Discrete
Mathematics \textbf{34}(2) (2020), 1011--1038, Theorem 3.1.
\url{https://doi.org/10.1137/19M1262206}.
\url{https://arxiv.org/abs/1905.05646}.
\bibitem{BB} [write] M. B\'ona and A. Burstein, \emph{Permutations with
exactly one copy of a monotone pattern of length $k$, and a generalization},
arXiv:2101.00332v3 (2021), Lemma~6.3 (with Theorems~4.1, 5.5 and~6.2).
\url{https://arxiv.org/abs/2101.00332}. Read 5~October 2026.
\bibitem{TSvol} [write] ProveIt, \emph{Transseries: the
polynomial-logarithmic calculus, series reversal at infinity, and
inversion},
\path{Analysis/Transseries/docs/series-and-transseries/Transseries_And_Inversion/transseries_and_inversion.tex},
\texttt{p0:def:three-inverses}, \texttt{p0:thm:lambert-core} and
\texttt{p0:thm:staircase}.
\end{thebibliography}
\end{document}
